\section{Identity-Indexed Local Tempo Routing}
\label{sec:method}

The counter begins after perception: it receives pose trajectories with supplied
identity indices and validity masks. Oracle tracks, tracks from a frozen common
frontend, and native method pipelines define different experimental protocols
and are never ranked together. Parameters are shared, whereas buffers, periodic
evidence, masks, and responses remain identity-indexed; no window combines
people. Figure~\ref{fig:framework} summarizes the proposed specification, whose
tensor layouts and learned modules remain \SyncRequired{collaborator source
binding}. For predicted tracks, an index denotes a current tracklet, not a
guaranteed person. Fragmentation, re-entry, switches, retirement, and any
reconciliation belong to the \SyncRequired{track lifecycle contract}.

\begin{figure*}[t]
  \centering
  \TempoRACFrameworkGraphic
  \caption{Conceptual framework reproduced in full from the project's editable
  framework deck. In the manuscript, $M_i$ denotes an identity-indexed counter
  instance with shared learned parameters and separate state, not a separately
  parameterized model. The training strip lists candidate objectives: tempo
  routing and track-continuity terms remain \texttt{SYNC-REQUIRED} until they
  are bound to frozen code. PAMS-TCC is attributed to PAMS
  \cite{gao2026pams}. The diagram is not implementation or result evidence.}
  \label{fig:framework}
\end{figure*}

\subsection{Masked identity-indexed representation}
For person $i\in\{1,\ldots,N\}$, let \(\bX_i\) contain $J$ joints with $C$
coordinates over $T_i$ frames, and let $m_i(t)$ mark a valid observation.
Overlapping supports \(\cW_{i\ell}\) select local windows. A shared encoder
\(\cE_\theta\) maps the masked samples to features, while the primary output
remains a person-wise count vector:
\begin{equation}
\begin{aligned}
 \bX_i &\in \Rset^{T_i\times J\times C}\qquad
 m_i\in\Bset^{T_i}\qquad
 \widehat{\bc}=[\hat c_i]_{i=1}^{N} \\
 \bz_{i\ell}(t) &=
 \cE_\theta\!\left(\bX_i[\cW_{i\ell}]\mid m_i[\cW_{i\ell}]\right)_t
 \quad t\in\cW_{i\ell}
\end{aligned}
\label{eq:representation}
\end{equation}
\noindent Here $\Bset$ denotes the binary alphabet. The shared $\theta$ does
not mix identity state. Window geometry, pose
normalization, padding, short tracks, and buffer lifecycle are
\SyncRequired{track tensor and lifecycle contract}. Masks accompany every
transformation, and windows never cross identity timelines. A scene total
$\sum_i\hat c_i$ is diagnostic. Permutation, isolated-mask, and short-track
tests will check these invariants once the source interface exists.

\subsection{Local masked periodic evidence}
Tempo must vary with the window rather than being fixed for the whole track.
Let $s_{i\ell}(t)$ denote the scalar evidence signal provisionally projected
from \(\bz_{i\ell}(t)\), and $q_{i\ell}(t)$ its windowed validity mask. We
define a mask-normalized non-circular autocorrelation and a normalized Fourier
score, then map both to a period estimate and confidence:
\begin{equation}
\begin{aligned}
 \bar s_{i\ell} &=
 \frac{\sum_t q_{i\ell}(t)s_{i\ell}(t)}
      {\sum_t q_{i\ell}(t)+\epsilon} \\
 A_{i\ell}(\tau) &=
 \frac{\sum_t q_{i\ell}(t)q_{i\ell}(t+\tau)
 [s_{i\ell}(t)-\bar s_{i\ell}]
 [s_{i\ell}(t+\tau)-\bar s_{i\ell}]}
 {\sum_t q_{i\ell}(t)q_{i\ell}(t+\tau)+\epsilon} \\
 F_{i\ell}(\omega) &=
 \frac{\left|\sum_t q_{i\ell}(t)[s_{i\ell}(t)-\bar s_{i\ell}]
 e^{-\mathrm{j}\omega t}\right|^2}
 {\max_{\omega'}\left|\sum_t q_{i\ell}(t)[s_{i\ell}(t)-\bar s_{i\ell}]
 e^{-\mathrm{j}\omega' t}\right|^2+\epsilon} \\
 \begin{bmatrix}\hat p_{i\ell}\\\gamma_{i\ell}\end{bmatrix} &=
 \Phi\!\left(\{A_{i\ell}(\tau)\}_{\tau\in\cT}\mathbin\Vert
              \{F_{i\ell}(\omega)\}_{\omega\in\Omega}\right)
\end{aligned}
\label{eq:period}
\end{equation}
\noindent Indices $(i,\ell)$ expose within-track drift. Projection, centering,
domains and units, interpolation, confidence, and fallback are
\SyncRequired{period estimator}; Eq.~\ref{eq:period} fixes only the semantics of
local masked evidence. Confidence must measure available evidence rather than a
post-routing label, and low-support windows must be masked or use one frozen
fallback in training and evaluation. Since periodic visibility can induce
spectral peaks, the implementation must set out-of-range masks and minimum
lag-pair support and pass periodic/nonperiodic missingness controls. It must also
choose pair-dependent centering or another irregular-sample estimator
\cite{lomb1976unequal}.

\subsection{Soft routing over shared response experts}
Let \(\cK=\{\mathrm{fast},\mathrm{medium},\mathrm{slow}\}\) index shared
experts. A proposed router consumes the window representation, local period,
and confidence, and returns simplex weights. Each expert emits a continuous
local response rather than an independently rounded count:
\begin{equation}
\begin{aligned}
 \bu_{i\ell} &=
 \rho_\psi([\bz_{i\ell}\mathbin\Vert\hat p_{i\ell}\mathbin\Vert\gamma_{i\ell}]) \\
 \bg_{i\ell} &= \softmax(\bu_{i\ell}/\tau_{\mathrm r}) \\
 g_{i\ell k}&\geq0\qquad \sum_{k\in\cK}g_{i\ell k}=1 \\
 \br^{(k)}_{i\ell} &=
 \cH^{(k)}_{\phi_k}(\bz_{i\ell})\qquad k\in\cK
\end{aligned}
\label{eq:routing}
\end{equation}
\noindent Soft weights retain uncertainty at tempo transitions. Router inputs,
temperature, expert architecture, gradients, sharing, and fallback remain
\SyncRequired{router and expert implementation}. Fast/medium/slow are intended
scales, not observed specializations: paired controls must test expert use,
collapse, frozen tempo bands, and hard/uniform routing. Shared experts keep
parameter count independent of the number of tracks, but no runtime or memory
claim is enabled without profiling.

\subsection{Response reconstruction and one track decode}
The final stage mixes expert responses before merging overlapping windows. Let
$a_{i\ell}(t)$ be a taper on window \(\ell\), and let
$q_{i\ell}(t)$ exclude invalid samples. Normalized overlap-add forms a single
track response $R_i$, after which one decoder extracts events:
\begin{equation}
\begin{aligned}
 \widetilde{r}_{i\ell}(t) &=
 \sum_{k\in\cK}g_{i\ell k}r^{(k)}_{i\ell}(t) \\
 R_i(t) &=
 \frac{\sum_{\substack{\ell\\t\in\cW_{i\ell}}}
 a_{i\ell}(t)q_{i\ell}(t)\widetilde{r}_{i\ell}(t)}
 {\sum_{\substack{\ell\\t\in\cW_{i\ell}}}
 a_{i\ell}(t)q_{i\ell}(t)+\epsilon} \\
 \widehat{\cE}_i &= \cD(R_i\mid m_i)\qquad
 \hat c_i=|\widehat{\cE}_i|
\end{aligned}
\label{eq:reconstruction}
\end{equation}
\noindent This order avoids counting every window independently; it does not
guarantee one peak. Decoded samples require aligned timestamps and coverage
$\sum_\ell a_{i\ell}(t)q_{i\ell}(t)>0$. Taper, gaps, event confidence, peak
parameters, and decoder semantics remain \SyncRequired{reconstruction and
decoder}. PAMS-TCC is the documented starting objective \cite{gao2026pams};
$\cL_{\mathrm{PAMS\text{-}TCC}}+\lambda_{\mathrm r}\cL_{\mathrm{route}}
+\lambda_{\mathrm t}\cL_{\mathrm{track}}$ is only a synchronization contract.
The latter terms, their existence, formulas, weights, samples, supervision, and
gradient paths must bind to code; absent terms will be deleted. Overlap-add
preserves continuous evidence before one masked trajectory-level decode. A
boundary test will place an event at several offsets and test one-time decoding.

The supervision audit must inspect loaders, pseudo-targets, auxiliary losses,
selection, stopping, calibration, and inference fallbacks at one commit. Until
it passes, the paper states only the intended absence of person-wise count and
period labels in the counting objective and discloses perception supervision.
